\section{Scaling equivalence between parameters and reasoning tokens}
\label{sec:scaling}

This section reports the fitted trade-off between model capacity and test-time reasoning.
The model is

\begin{equation}
\operatorname{logit}\bigl(\mathrm{accuracy}\bigr) = \beta_P \log_2\bigl(\mathrm{params}_B\bigr)
+ \beta_R \log_2\bigl(1 + \mathrm{reasoning\ tokens}\bigr) + \mathrm{controls},
\label{eq:scaling}
\end{equation}

with controls for topology, context sharing, prompt level and a 32B context-cap indicator.
A categorical companion replaces the continuous reasoning term with rung dummies. Both are
fitted per benchmark, and both are fitted twice: once on the full core view with the cap
indicator, and once on a subset that drops the 32B cells at b8192 and unlimited entirely.
These are descriptive fits over a partially complete sweep, not asymptotic scaling laws.

\subsection{Fitted coefficients}

Table~\ref{tab:scaling-coefficients} gives all 40 fits, and
Figure~\ref{fig:scaling-coefficients} displays the pooled and per-topology coefficients.

\intable{tab_scaling_coefficients}

\begin{figure}[H]
\centering
\includegraphics[width=\linewidth,height=0.34\textheight,keepaspectratio]{fig_scaling_coefficients.pdf}
\caption{Capacity and reasoning coefficients by fit group and benchmark, with 95 percent
intervals. The per-topology fits are reported here for the first time; the source memo
printed only the pooled row.}
\label{fig:scaling-coefficients}
\end{figure}

On the pooled fit, $\beta_P$ is 0.230 on GPQA, 0.384 on MMLU-Pro, 0.399 on TruthfulQA and
0.275 on MATH, and $\beta_R$ is 0.032, 0.057, 0.042 and 0.032. Deviance $R^2$ is 0.769,
0.866, 0.906 and 0.665. The mathematical benchmark is the least well explained by this
specification, which is consistent with its accuracy surface being non-monotone in capacity.

The ratio $\beta_P/\beta_R$ is the number of reasoning-token doublings that substitute for
one parameter doubling: 7.22 on GPQA, 6.72 on MMLU-Pro, 9.50 on TruthfulQA and 8.67 on
MATH. Read as a local marginal rate, one doubling of realized reasoning tokens buys only
about 0.10 to 0.15 of a parameter doubling. Since a model cannot realistically spend seven
to nine doublings of thinking tokens, the continuous form understates what the reasoning
axis actually delivers, because most of its value is in the discrete step from disabled to
enabled thinking rather than in the marginal token.

The per-topology fits, which the source analysis computed but did not report, show that
$\beta_P$ increases with coordination. On GPQA it is 0.203 for single agent, 0.216 for
independent and 0.245 for decentralized. Larger models extract more from being placed in a
debate than smaller models do, which is a capacity-by-topology interaction that the pooled
fit averages away.

\subsection{Reasoning rungs as parameter multipliers}

Table~\ref{tab:rung-equivalents} in Appendix~\ref{app:fits} converts each rung's logit
gain into an equivalent parameter multiplier by dividing by $\beta_P$ from the same fit.

For the pooled fit the multipliers against thinking disabled are:

\begin{center}
\small
\begin{tabular}{lrrrr}
\toprule
Benchmark & b512 & b2048 & b8192 & unlimited \\
\midrule
GPQA & 1.89 & 3.04 & 4.01 & 3.92 \\
MMLU-Pro & 1.99 & 3.23 & 3.85 & 3.90 \\
TruthfulQA & 1.93 & 2.14 & 2.27 & 2.34 \\
MATH & 0.71 & 2.11 & 3.97 & 3.50 \\
\bottomrule
\end{tabular}
\end{center}

Enabling thinking at b2048 is worth roughly two to three times the parameters, and b8192 is
worth roughly 2.3 to 4.0 times, depending on the task. The unlimited rung does not
materially exceed b8192 on any benchmark and is lower on two.

The mathematical benchmark has the one anomaly: b512 is worth 0.71 times the parameters,
that is worse than disabling thinking. A short thinking budget on a task requiring extended
derivation appears to be actively harmful, presumably because the model begins a derivation
it cannot finish within the budget and then answers from an incomplete state. This is the
only place in the sweep where a reasoning rung is worse than no reasoning at all.

Figure~\ref{fig:size-reasoning-heatmap} shows the underlying accuracy surface with the
context-capped cells marked.

\begin{figure}[H]
\centering
\includegraphics[width=\linewidth,height=0.30\textheight,keepaspectratio]{fig_size_reasoning_heatmap.pdf}
\caption{Accuracy over the capacity and reasoning grid. Red hatching marks cells carrying
the 32B context-cap flag, which is where the terminal capacity contrast is least
trustworthy.}
\label{fig:size-reasoning-heatmap}
\end{figure}

The surface is not a uniform vertical shift. GPQA at 0.6B stays near 0.26 to 0.28
regardless of reasoning budget, while 14B and 32B with a large budget reach roughly 0.60 to
0.63. MMLU-Pro has the smoothest surface, reaching 0.84 to 0.86 at the top. TruthfulQA
saturates near 0.81 to 0.82 for large models with reasoning. MATH is the least monotone,
with several small and mid-capacity models dipping at b512 and the best region at 8B or 14B
with b8192 rather than at 32B.

Reasoning is therefore most valuable once the base model has enough capacity to exploit it.
For the weakest models on the hardest benchmark, additional thinking alone does very little.
Table~\ref{tab:size-reasoning-grid} in Appendix~\ref{app:fits} gives the complete grid of
accuracy, realized token counts and cost proxy for the pooled fit and for each topology
separately, 598 rows.

\subsection{Plateaus}

Table~\ref{tab:plateaus} gives the adjacent-rung contrasts on both scaling axes with a
plateau flag.

\intable{tab_plateaus}

The reasoning axis plateaus unambiguously. The unlimited against b8192 contrast is $-0.005$
$[-0.011, 0.001]$ on GPQA, $+0.001$ $[-0.003, 0.005]$ on MMLU-Pro, $-0.000$ $[-0.004,
0.003]$ on TruthfulQA and $-0.008$ $[-0.013, -0.003]$ on MATH. Three intervals contain zero
and the fourth is negative. Removing the token budget entirely does not beat an 8{,}192-token
budget.

The b8192 against b2048 contrast is still positive on three benchmarks, at $+0.020$
$[0.014, 0.027]$, $+0.016$ $[0.011, 0.022]$ and $+0.039$ $[0.031, 0.047]$, and zero on
TruthfulQA at $+0.000$ $[-0.004, 0.004]$. The reasoning axis therefore saturates between
b8192 and unlimited, not before.

The capacity axis appears to plateau between 14B and 32B, but this conclusion is weaker.
The 32B block at b8192 and unlimited holds only 32 and 30 design cells against roughly 47
elsewhere, and those cells are context-capped. The plateau is measured exactly where the
data are thinnest and most confounded. It cannot be separated from the cap with the present
sweep, and Section~\ref{sec:limitations} states this as an outstanding item.

\subsection{Nearest empirical equivalences}

Table~\ref{tab:nearest-equivalences} in Appendix~\ref{app:fits} lists, for each
low-reasoning target configuration,
the smaller model with a larger reasoning budget whose accuracy is closest, together with
the ratio of their cost proxies. Figure~\ref{fig:equivalence-arrows} draws the same
substitutions on the cost-accuracy surface.

\begin{figure}[H]
\centering
\includegraphics[width=\linewidth,height=0.52\textheight,keepaspectratio]{fig_equivalence_arrows.pdf}
\caption{Cost-accuracy surface by benchmark, with an arrow from each low-reasoning target
configuration to its nearest cheaper equivalent. Arrows pointing left and level identify
substitutions that preserve accuracy at lower inference cost.}
\label{fig:equivalence-arrows}
\end{figure}

The numbers quoted below are from the pooled \emph{all systems} scope, matching the fits
above. The table also carries the four per-topology scopes, and the substitutions they
identify differ: the single-agent scope, for example, puts 32B with thinking disabled at
0.372 on GPQA rather than 0.415, because it excludes the accuracy that coordination adds.
Any equivalence claim is therefore scope-specific and should be read with its scope stated.

Substituting reasoning for parameters is strongly favourable at the top of the capacity
ladder. Against 32B with thinking disabled, 4B at b512 reaches 0.426 on GPQA against 0.415,
at 0.12 times the inference cost; 8B at b512 reaches 0.701 on MMLU-Pro against 0.703, at
0.30 times the cost; 4B at b2048 reaches 0.758 on MATH against 0.741, at 0.24 times the
cost. On TruthfulQA the nearest candidate falls short, with 4B at b2048 reaching 0.680
against 0.721, though at 0.19 times the cost.

The substitutions are unfavourable at the bottom of the ladder, which bounds the claim.
Against 4B with thinking disabled, 1.7B at b8192 reaches 0.325 on GPQA against 0.397 at
1.53 times the cost, 0.558 on MMLU-Pro against 0.598 at 1.25 times the cost, and 0.674 on
MATH against 0.711 at 1.06 times the cost. In all three the smaller model is both less
accurate and more expensive.

Thinking does not rescue a model that is too small for the task, and it can cost more than
the parameters it was meant to replace, because the cost proxy multiplies parameters by
tokens and a long thinking trace on a small model exceeds a short answer from a large one.
The pattern across the whole table is that substituting reasoning for parameters is
favourable when the smaller model is within roughly one capacity step of the target and
unfavourable at two steps or more.
